% ПРИЛОЖЕНИЕ, без номер на раздел.
% Изходният текст НЕ се вмъква тук. Хранилището е публично и една таблица с
% указател към файловете е по-полезна от двадесет страници преписан код.
\section*{ПРИЛОЖЕНИЕ}
\label{sec:appendix}
\addcontentsline{toc}{section}{ПРИЛОЖЕНИЕ}

Целият изходен текст, вложеният сценарий, тестовете, уеб клиентът и пълният
изход на командите \code{quorum demo} и \code{quorum bench} са в публичното
хранилище \url{https://github.com/Alex-Tsvetanov/quorum}. Построяване и
изпълнение, без външни зависимости, само с компилатор за C++20 и CMake:
\code{cmake -S . -B build -DCMAKE\_BUILD\_TYPE=Release}, след това
\code{cmake --build build} и \code{./build/quorum demo}. Уеб страницата се вдига с
\code{./build/quorum serve}. Указателят кое изчисление в кой файл живее е в
Табл.~\ref{tab:sources}; броят редове на ядрото е от \code{wc -l} към състоянието,
от което са получени числата в Раздел~\ref{sec:results}, а редовете на
представянето са към състоянието с уеб клиента.

\begin{table}[H]
\centering
\caption{Изходен текст: къде живее всяко изчисление}
\label{tab:sources}
\small
\begin{tabular}{@{}L{3.2cm}L{7.4cm}r@{}}
\toprule
\textbf{Файл в \code{src/}} & \textbf{Съдържание} & \textbf{Редове} \\
\midrule
\code{json.cpp} & четец и записвач на JSON с коментари & 368 \\
\code{model.cpp} & сценарий, инварианти, изведени величини & 271 \\
\code{cpm.cpp} & топологично подреждане, критичен път, PERT & 245 \\
\code{mcdm.cpp} & тегла по AHP, съгласуваност, TOPSIS & 244 \\
\code{portfolio.cpp} & разклоняване и ограничаване, лаком алгоритъм & 345 \\
\code{risk.cpp} & Монте Карло, триъгълно разпределение & 162 \\
\code{sensitivity.cpp} & сканиране по тегла и по бюджет & 132 \\
\code{report.cpp} & таблици и обяснения на отхвърлянето & 215 \\
\code{analyse.cpp} & общият конвейер към доклада и към JSON & 213 \\
\code{http.cpp}, \code{net.cpp}, \code{web.cpp} & HTTP/1.1, сокети, маршрутизация & 580 \\
\code{bench.cpp}, \code{main.cpp} & измерванията, командният ред и \code{serve} & 407 \\
\midrule
\code{include/quorum/} & дванадесет публични заглавни файла & 803 \\
\code{tests/} & 83 случая в девет групи и рамката за тестване & 1432 \\
\code{web/} & страницата: HTML, CSS и JavaScript без пакетна система & 538 \\
\code{scenarios/} & съставният сценарий с дванадесет проекта & 206 \\
\bottomrule
\end{tabular}
\end{table}

\textbf{Схема на сценария.} Сценарият е един JSON документ. Задължителни са
\code{budget}, \code{criteria} и \code{projects}. Критерият носи \code{id} и по
избор \code{direction}, \code{max} или \code{min}; теглата се задават пряко чрез
\code{weight} или се извеждат от квадратната положителна реципрочна матрица
\code{comparisons}. Ресурсът носи \code{id} и \code{capacity}. Проектът носи
\code{id}, оценка по всеки критерий в \code{scores}, по избор \code{requires} и
\code{excludes}, и непразен списък \code{activities}. Дейността носи \code{id},
\code{depends\_on} и \code{demand} по ресурс. Продължителността и стойността се
задават или с едно число, \code{duration} и \code{cost}, или с тройна оценка,
\code{optimistic}, \code{likely} и \code{pessimistic}, със същите представки за
стойността; едното число е случаят, в който трите точки съвпадат, затова за него
няма отделен път в кода.
